% working_intro.tex — restructured Nature-style intro draft (8-paragraph arc).
% Paragraphs marked NEW are new material; others are carried over from intro.tex
% with edits noted. Related Work section is untouched and remains in intro.tex.
% Search for PLACEHOLDER to find citation keys that need real .bib entries.

% ¶1 — What cognitive models are; the inverse payoff planted immediately.
% (intro.tex lines 1–3, with a new third sentence and "inverted" added to the last.)
Cognitive models cast psychological theories in probabilistic form, specifying how unobservable mental quantities (e.g., processing speed) give rise to observable behavior (e.g., reaction times).
These models are typically \textit{generative}; that is, they define a recipe for simulating realistic behavior from any given configuration of the underlying constructs.
Their scientific value, however, is realized in reverse: by inverting the model against observed behavior, latent constructs become measurable, so that a drift rate can quantify processing speed and a decision threshold can quantify response caution.
Traditionally, this machinery is handcrafted anew for each experiment: one model specification per task, simulated, inverted, and analyzed on its own.

% ¶2 — The forward meta-model moment. (intro.tex lines 4–6, unchanged.)
A recent development breaks with this convention by moving one level up, from models to \textit{meta-models}: single systems trained over an entire space of tasks or model specifications, which then generalize across that space.
The most prominent example to date is Centaur \citep{binz2025centaur}, a language model fine-tuned on trial-by-trial data from hundreds of psychological experiments that predicts human behavior in held-out tasks with remarkable fidelity.
Viewed through this lens, Centaur is a meta-model of the \textit{forward} direction: instead of implementing any particular cognitive model, it subsumes the behavior-generating machinery of many task-specific models into one system.

% ¶3 — NEW: the controversy as hinge; the asymmetry. (Replaces intro.tex lines 8–13.)
Yet the scientific payoff of cognitive modeling lies almost entirely in the opposite, \textit{inverse} direction: given observed behavior, what can we conclude about the latent constructs that produced it?
When a patient responds slowly, is it because evidence accumulates sluggishly or because they are cautious?
Answering such questions
 is what cognitive models are \textit{for}---measurement, theory testing, and computational phenotyping all operate in the inverse direction.
It is telling that the debate Centaur has ignited among cognitive scientists turns on precisely this point: critics object that a system predicting behavior without exposing interpretable mechanisms leaves the explanatory questions of psychology untouched \citep{PLACEHOLDER-centaur-critiques}.
We take no side in this debate here; we instead read it as a diagnosis.
The meta-modeling revolution has so far arrived only in the forward direction, where its scientific yield is genuinely contested, and not in the inverse direction, where its yield would be immediate.
Centaur itself offers no help: because it does not preserve information about latent constructs, it cannot tell a researcher what an observed pattern of responses implies about processing speed, response caution, or memory capacity.
Its inverse counterpart---a meta-model that maps behavior to posterior distributions over the parameters of interpretable cognitive models, across an entire space of such models---does not yet exist.

% ¶4 — NEW: the workflow-pain paragraph, from the researcher's chair.
Consider what inverse inference actually looks like in practice.
A researcher modeling response times rarely fits one model; they fit a lattice of them.
Does the experimental manipulation shift the drift rate, the decision threshold, or both \citep{voss2019sequential, von2022mental}?
Each hypothesis is a model variant, and each answer requires fitting the variant that embodies it against the variants that do not.
Is across-trial variability in the starting point needed?
Two more fits.
Should the threshold be allowed to collapse over time?
Two more \citep{evans2019assessing}.
Because the posterior $p(\rvtheta \mid \rvy)$ is rarely tractable for such models, each node of this lattice is visited with iterative Monte Carlo methods whose runtimes range from hours to days \citep{stricklandRacingRemember2018, mileticParameterRecovery2017}, and whose convergence must be diagnosed and often coaxed by hand.
A hierarchical extension, an added covariate, or a reviewer's request for one more comparison restarts the process.
The result is a quiet mismatch at the heart of the field: theory testing is conducted over \textit{spaces} of models, while our inference machinery is built, and paid for, one model at a time.

% ¶5 — SBI and amortization. (intro.tex lines 15–17, with a plain-language gloss
% of amortization added to the second sentence.)
The natural starting point for addressing this mismatch is neural simulation-based inference \citep[SBI;][]{cranmer2020frontier, deistler2025simulation}, which employs neural networks trained on simulated data to avoid iterative estimation.
In particular, \textit{amortized methods} \citep{burkner2023some} learn from simulated parameter–observation pairs $(\rvtheta, \rvy)$: the cost of estimation is paid once, up front and in simulation, after which the trained network samples from the posterior $p(\rvtheta \mid \rvy^*)$ for any new observation $\rvy^*$ in seconds, without re-estimation.
Cognitive models are, in many ways, ideal generators of the required synthetic data: they are easy to implement, fast to simulate (often in the millisecond range), and low-dimensional.
Accordingly, SBI applications in cognitive modeling abound \citep[e.g.,][]{radev2020amortized, wieschen2020jumping, fengler2021likelihood, boelts2022flexible, von2022mental, ghaderi2023general, sokratous2023ask, schumacher2023superstat, schumacher2024validation, rmus2024artificial, hato2025, schaefer2025conflictModels, scholten2026diffusionLens, kvam2024using, kvam2025comparing, belov2026supervised, wu2026testing}.

% ¶6 — NEW: the granularity argument; meta-amortization defined; taxonomy preview.
% (Replaces intro.tex lines 18–26.)
Standard amortization, however, inherits the very granularity that creates the mismatch.
The network is trained for a single, fully specified model, so its validity ends exactly where the specification does: fix a parameter, add a covariate, or swap a boundary function, and the trained network is obsolete, its up-front cost re-paid at every node of the lattice.
An inference tool matched to the scope of the scientific question must therefore be amortized at the level of the family, or beyond it, rather than the individual specification.
We call this \textit{meta-amortization}: paying the simulation and training cost once for an entire space of model specifications, after which any dataset can be inverted under any model in that space, instantly and without retraining.
To make the scope of such claims precise, we organize model spaces into a taxonomy of \textit{species} (a fully specified model), \textit{families} (species related by parameter constraints), and \textit{classes} (families sharing a computational principle, such as sequential sampling), and we demonstrate meta-amortization at both the family and the class level.

% ¶7 — The contribution. (intro.tex lines 28–32, with the inference-quality claim
% rephrased in the field's vocabulary — check wording against what the results support.)
The current work delivers such a meta-amortized inference engine, building on a stream in SBI that extends amortization across models \citep{chang2025amortized, gloeckler2024all, elsemuller2023sensitivity, schroder2023simultaneous}:
training a single inference network that remains valid across many, potentially infinitely many, structurally similar models.
We dub our architecture the \textbf{CogFormer}, directly inspired by the \textbf{SimFormer} \citep{gloeckler2024all}.
CogFormer adapts to any configuration within a model collection's combinatorial parameter space and infers its components accordingly; in parameter recovery and calibration, it matches networks amortized for a single model, so the gain in scope costs little in the currency the field actually audits.
In this work, we taxonomize the scope of amortization in cognitive modeling and demonstrate compelling results from the first meta-amortized inference engine for cognitive modeling: where Centaur offers a meta-model of how behavior is generated, CogFormer offers its inverse---a meta-model of how latent cognition is recovered from behavior.

% ¶8 — NEW: the shared-instrument significance paragraph.
% (Final sentence trimmed to avoid duplicating the ¶7 capstone.)
Meta-amortization also changes who pays the cost of inference, and when.
Under per-model amortization, every lab pays for its own network, for every model, before the first dataset is fit---an economy that has largely confined amortized inference to methodologically specialized groups.
A meta-amortized estimator inverts this economy: the training cost is paid once, by whoever trains the network, and what the community receives is an instrument.
A researcher need only supply data and a model specification to obtain calibrated posteriors in seconds, for whichever variant of the family their theory demands.
The precedent is instructive: hierarchical modeling toolboxes spread through cognitive science not because they introduced new models, but because they removed the implementation burden of fitting existing ones \citep{PLACEHOLDER-hddm-emc2-toolboxes}.
An inference engine amortized over a model class removes the analogous burden for inference itself, and it removes it at a scale that matters beyond the laboratory: computational phenotyping, which seeks model parameters as clinical measurements in large and heterogeneous samples \citep{PLACEHOLDER-huys-phenotyping}, is precisely the setting in which refitting per person, per model variant, does not scale.
The result is not a rival to interpretable cognitive models, but the machinery that finally lets them be used at the scale the forward direction has already reached.
